\section{Discussion}
\label{sec:discussion}

We interpret the results of Section~\ref{sec:results} against the four research questions of Section~\ref{sec:introduction}, derive three design guidelines for practitioners deploying multi-SLM software-engineering pipelines, and characterize the threats to validity that scope our claims in Section~\ref{sec:limitations}.

\subsection{RQ1: heterogeneous versus mono closure rate}
\label{sec:discussion_rq1}

\textbf{RQ1 answer.} Heterogeneous cells match rather than dominate the strongest mono baseline on aggregate closure rate, and \emph{substantially} exceed it when decomposed by mutation operator (Path~1) or by stage attribution.

On Path~1 metric mean $\bar{Y}$ (Eq.~\ref{eq:metric_mean}), the best hetero cell H6 ($\bar{Y}_{H6,1} = 0.882$) does not statistically differ from the best mono cell M3 ($\bar{Y}_{M3,1} = 0.900$; Tukey $p_{\textsf{adj}} = 0.999$). The naive interpretation would be that heterogeneity offers no benefit. This interpretation is unstable across two orthogonal decompositions:

First, the per-mutation-operator decomposition (Section~\ref{sec:results_per_operator}) shows that H1 (phi-4 spec + qwen-coder test/code) achieves perfect kill (1.000) on comparison and negate\_bool operators and $0.939$ on arithmetic --- all three exceeding the best mono baseline by $\geq 0.09$ percentage points. Any single mono cell wins on some operators and loses on others; the H1 composition achieves the operator-wise maximum on four of five operator families simultaneously. \emph{The mono ceiling on any given operator is not a fundamental capability ceiling.}

Second, the stage-attribution decomposition (Section~\ref{sec:results_stage_attribution}), evaluated on the strict-AND valid-closure fraction $V$ (Eq.~\ref{eq:valid_closure_frac}, denominator = total measurement rows), shows that H11 and H4 exceed the strongest mono ceiling on Path~1 outright ($V_{H11,1} = 0.813$ and $V_{H4,1} = 0.800$ vs.\ $V_{M4,1} = 0.773$). This is a $V$-scale claim; on the metric-mean $\bar{Y}$ scale, M3 still leads. Both hetero cells specialize at the code and test stages while permitting a weaker model at spec (H4 uses llama3.2 at code, H11 uses phi-4 at spec) --- validating that the heterogeneity benefit compounds when specialization matches the demanded stage strength.

The RQ1 answer therefore has two parts: heterogeneity does not improve \emph{aggregate} metric-mean closure over the strongest mono baseline, but it does improve the \emph{strict-AND valid-closure fraction} for two hetero cells and the \emph{per-operator maximum} for one hetero cell (H1) in ways that any single mono baseline cannot achieve.

\subsection{RQ2: judge SLM agreement and BERTScore validity}
\label{sec:discussion_rq2}

\textbf{RQ2 answer.} The external judge SLM correlates with mutation kill rate ($r_1 = 0.192$, weak but significant) and reference pass rate ($r_2 = 0.523$, moderate), and correlates only weakly per benchmark with BERTScore ($r_3^{\textsf{HE}} = 0.153$, $r_3^{\textsf{MBPP}} = 0.212$; both $r^2 \leq 0.045$). The aggregate $r_3 = 0.022$ ($p = 0.359$) is a Simpson's-paradox attenuation from pooling; the per-benchmark weak-but-significant correlations have a structural directional signature that flips across benchmarks.

\begin{table}[!htbp]
\centering
\small
\caption{Judge-metric agreement decomposition by path. Categories from Algorithm~2's decision branches. \textbf{Bold}: Path~3's Metric~$\checkmark$/Judge~$\times$ share ($24.5\%$) is more than double Path~1's ($11.3\%$), consistent with BERTScore functioning as a surface-form similarity rather than a semantic-equivalence signal.}
\label{tab:disagreement_by_path}
\begin{tabular}{lrrrrr}
\toprule
Path & Both $\checkmark$ & Both $\times$ & Metric $\checkmark$/Judge $\times$ & Metric $\times$/Judge $\checkmark$ & N/A \\
\midrule
Path 1 & 1,142 (58.1\%) & 26 (1.3\%) & 222 (11.3\%) & 51 (2.6\%) & 523 (26.6\%) \\
Path 2 & 978 (49.8\%) & 408 (20.8\%) & 83 (4.2\%) & 370 (18.8\%) & 125 (6.4\%) \\
Path 3 & 1,076 (54.8\%) & 62 (3.2\%) & \textbf{480 (24.5\%)} & 173 (8.8\%) & 171 (8.7\%) \\
\bottomrule
\end{tabular}
\end{table}

Stratifying the Path~3 rows by Algorithm~\ref{alg:validity} decision reason (Table~\ref{tab:disagreement_by_path}) shows that $24.5\%$ of all Path~3 rows are \textsf{false\_closure\_candidate}: the BERTScore threshold is met but the judge rating is below $\rho$. A further $8.8\%$ are \textsf{metric\_false\_negative}: the judge accepts the docstring pair but the BERTScore threshold is not met. Path~3 exhibits the highest disagreement of any path, more than double Path~1's rate ($11.3\%$ false-closure candidates).

Table~\ref{tab:disagreement_p3_by_benchmark} decomposes this disagreement by source benchmark.

\begin{table}[htbp]
\centering
\small
\caption{Path~3 judge--metric disagreement by source benchmark. \textbf{The disagreement flips direction across benchmarks}: MBPP over-credits equivalence (BERTScore high but judge disagrees), HumanEval under-credits (judge accepts but BERTScore does not).}
\label{tab:disagreement_p3_by_benchmark}
\begin{tabular}{lrrr}
\toprule
Benchmark & Both agree valid & \textsf{false\_closure\_candidate} & \textsf{metric\_false\_negative} \\
\midrule
HumanEval ($n = 550$) & $39.6\%$ & $12.2\%$ & $\mathbf{30.0\%}$ \\
MBPP ($n = 1{,}412$) & $60.8\%$ & $\mathbf{29.2\%}$ & $0.6\%$ \\
\bottomrule
\end{tabular}
\end{table}

\textbf{The disagreement flips direction across benchmarks.} On MBPP, BERTScore over-credits equivalence in $32.0\%$ of rows (versus under-crediting in only $0.6\%$); on HumanEval, BERTScore under-credits in $32.9\%$ of rows (versus over-crediting in only $13.4\%$). A $2 \times 2$ chi-square test on the (benchmark, disagreement type) contingency yields $\chi^2 = 368.0$ ($p < 0.0001$, $\text{df} = 1$), so the flip in disagreement direction is statistically categorical, not marginal. This structural signature is precisely what a surface-form similarity signal should produce and a semantic-equivalence signal should not: MBPP docstrings share a compact task-description vocabulary (``\texttt{return $X$ of $Y$}'', ``\texttt{check if $Z$ is $P$}''), so BERTScore assigns high F1 even when the described behaviors diverge; HumanEval docstrings vary in stylistic detail (examples, doctests, hints, edge-case notes), so BERTScore penalises stylistic drift even when the core semantics are preserved.

Combined with the very weak per-benchmark correlations ($r_3^{\textsf{HE}} = 0.153$, $r_3^{\textsf{MBPP}} = 0.212$; both $r^2 \leq 0.045$; see Section~\ref{sec:results_contamination}), this decomposition supports a stronger conclusion than ``BERTScore is a poor closure metric on Path~3''. It supports the specific claim that \textbf{BERTScore is functioning as a surface-form similarity signal rather than a semantic-equivalence signal on the docstring--docstring closure task}. This finding scopes our recommendation that judge-SLM adjudication is not an optional supplement to automated closure metrics on Path~3, but a necessary component of the validity decision.

\subsection{RQ3: per-stage bottleneck and the phi-4-in-test-stage pathology}
\label{sec:discussion_rq3}

\textbf{RQ3 answer.} Two orthogonal analyses converge on complementary answers. Under the min-baseline attribution rule (Section~\ref{sec:results_stage_attribution}), the spec stage is the plurality bottleneck (15 of 33 hetero (cell, path) triples), followed by the test stage (10) and the code stage (8). Under the mixed-effects marginal-coefficient rule (Section~\ref{sec:results_mixed_logit}, Table~\ref{tab:mixed_logit}), test-stage substitutions carry every significant per-stage attribution ($p < 0.01$ for llama3.2, mistral-small, phi-4 at $L_{\textsf{test}}$; all spec- and code-stage coefficients null at $\alpha = 0.05$). Both results are true and describe different aspects of the same underlying structure: the spec stage is most often the weakest link when a cell is bottlenecked, but the \emph{marginal} effect of substituting one model for another at a given stage is only significant at the test stage. The design implication (G2, Section~\ref{sec:discussion_capability}) that stage-model choice at the test stage matters most is supported by the mixed-effects result; the spec-stage plurality-bottleneck finding motivates the complementary G1 recommendation to place the strongest model at the bottleneck stage regardless of which stage that is for a given cell. A specific per-model pathology --- phi-4 at $L_{\textsf{test}}$ --- accounts for the largest single-cell attribution to the test stage and is the mechanism behind the mixed-effects test-stage coefficient.

Three cells assign phi-4 to $L_{\textsf{test}}$: M2 (phi-4 mono), H2 (qwen3.6 spec, phi-4 test, qwen-coder code), and H10 (mistral spec, phi-4 test, qwen-coder code). Path~1 \textsf{structural\_NA} rates for these cells are $44\%$ (M2), $45\%$ (H2), and $55\%$ (H10), corresponding to \textsf{filter\_reason} = \textsf{all\_dropped} on approximately half of all attempted samples. Cells that use phi-4 at other stages (H1 with phi-4 at $L_{\textsf{spec}}$ = 9\%; H6 with phi-4 at $L_{\textsf{code}}$ = 17\%; H11 with phi-4 at $L_{\textsf{spec}}$ = 5\%) exhibit substantially lower NA rates. \emph{The pathology is stage-specific}: phi-4 as $L_{\textsf{spec}}$ or $L_{\textsf{code}}$ works well; phi-4 as $L_{\textsf{test}}$ produces tests that fail the validity-filter pre-condition on roughly half of samples.

The mechanism, at operator resolution, is that phi-4-generated tests concentrate their failure on arithmetic and boundary mutations --- exactly the operators where phi-4 mono itself is competitive. H10's arithmetic kill rate ($0.351$) and boundary kill rate ($0.449$) are both substantially below the mono ceiling ($0.806$ and $0.748$ respectively). This is not a case of a weak model producing weak tests; it is a case of a model whose test-generation output is systematically misaligned with the defect families it should be strongest at detecting.

The practical implication for pipeline design is that per-model capability, as measured on a single-artifact task (docstring generation, or code synthesis), does not directly predict per-model capability at generating tests for that same artifact. The phi-4-in-test-stage effect observed here --- competitive at single-stage self-generation, degraded when generating tests for an upstream-produced docstring --- is the cross-stage interaction that RQ1--RQ3 collectively address.

\subsection{RQ4: false-closure rate}
\label{sec:discussion_rq4}

\textbf{RQ4 answer.} At the strong-success bar defined below, the mean false-closure rate across the 37 (cell, path) groups is $11.4\%$; the maximum non-null-cell rate is $28.6\%$ (H10 on Path~1); the null cell N1 (word-shuffled input) exhibits a rate of $26.3\%$, which serves as an empirical noise floor.

\input{../tables/tab_false_closure.tex}

The false-closure rate is the fraction of \emph{metric-success} rows within a (cell, path) group whose judge rating is below $\rho$. Metric success here uses per-path \emph{strong-success} thresholds --- kill rate $K \geq 0.70$ (Path~1), full reference pass $R = 1.00$ (Path~2), BERTScore $B \geq 0.80$ (Path~3), as implemented in \texttt{analyze/load\_results.py} --- deliberately stricter than the validity threshold $\tau = 0$ of Eq.~\ref{eq:valid1}--\ref{eq:valid3}. The strict bar isolates rows where the metric \emph{confidently} claims success, which is the population on which a judge veto is most diagnostic; at the permissive $\tau = 0$ bar, marginal metric signals (e.g.\ a single killed mutant) would dominate the count. This analysis therefore does \emph{not} equal the \textsf{false\_closure\_candidate} rate of Algorithm~\ref{alg:validity}, which is evaluated at $\tau = 0$. Table~\ref{tab:false_closure} reports the highest-false-closure (cell, path) triples at the strong-success bar.

The N1 rate of $26.3\%$ on Path~1 is important as a metric-soundness diagnostic. If the automated metric were a perfect closure signal, we would expect N1 (which feeds word-shuffled code into $L_{\textsf{spec}}$) to score near zero on the metric. It does not: $\sim$$26\%$ of N1 rows still exhibit metric $> \tau$. This reflects the fact that the mutation kill rate on filtered tests, even generated from shuffled code, retains a $\sim$$26\%$ false-positive floor. We treat this as the empirical noise floor of the metric under our threshold policy.

Under this interpretation:

\begin{itemize}
  \item Non-null cells with false-closure rates below the N1 floor ($< 26\%$) exhibit false-closure signals that are at or below noise. Every hetero cell except H10 lies below this floor on Path~1.
  \item H10 ($28.6\%$ on Path~1) is the only non-null cell above the noise floor, consistent with the phi-4-in-test-stage pathology of Section~\ref{sec:discussion_rq3}.
  \item Mono cells cluster in the $[10\%, 20\%]$ range, well below the noise floor.
\end{itemize}

The false-closure result is therefore not a limitation but a \emph{characterization} of the metric's operating regime, and it is bar-specific in one important respect. At the permissive $\tau = 0$ bar (any metric signal counts as success), rates rise everywhere --- N1 reaches $43.9\%$ and H10 $29.0\%$ on Path~1 --- so the \emph{noise floor} rises faster than the cell rates and the ``only H10 exceeds the floor'' statement holds only at the strong-success bar. What is threshold-robust is the ordering: H10 carries the highest non-null false-closure rate on Path~1 at both bars ($28.6\%$ strong; $29.0\%$ at $\tau = 0$, versus $22.7\%$ for the next non-null cell), so the phi-4-in-test-stage attribution does not depend on the bar choice.

\paragraph{Per-cell judge-leniency check.} A residual concern is whether between-cell contrasts under strict-AND are confounded by the judge treating some cells' outputs more leniently than others (verbosity or stylistic bias). To test this, Table~\ref{tab:judge_leniency} reports the per-cell mean judge rating restricted to rows where the metric already indicates high quality ($Y > 0.67$). If the judge were style-neutral, restricting to high-metric rows should yield near-constant mean judge ratings across cells. Empirically, the mean judge rating on high-metric rows ranges $[2.955, 3.317]$ across the 17 non-null cells (standard deviation $0.102$ on the 0--4 scale). This $\sim$$0.3$-point range across cells is small (residual after controlling for metric) but non-zero: the strongest leniency contrast is H7 ($3.317$) vs.\ M1 ($2.955$), a $0.36$-point difference. The between-cell strict-AND contrasts reported throughout the paper should therefore be interpreted with a $\sim$$5$~pp additive uncertainty attributable to residual judge leniency drift.

\begin{table}[t]
\centering
\small
\caption{Per-cell mean judge rating, overall and restricted to high-metric rows ($Y > 0.67$). Restricting to high-metric rows controls for artifact quality: any residual per-cell variation in mean judge rating is attributable to the judge treating some cells' outputs more leniently than others (e.g., verbosity-of-output effect). Standard deviations of the high-metric-only column across cells within a stratum quantify the judge-leniency confound.}
\label{tab:judge_leniency}
\begin{tabular}{@{}lrrrr@{}}
\toprule
Cell & $n$ & mean $J$ (all) & $n_{\text{high-Y}}$ & mean $J$ (high-Y only) \\
\midrule
H1 & 258 & 2.891 & 125 & 3.256 \\
H10 & 184 & 2.712 & 67 & 3.030 \\
H11 & 221 & 2.959 & 109 & 3.248 \\
H2 & 187 & 2.781 & 77 & 3.091 \\
H3 & 185 & 2.795 & 90 & 3.122 \\
H4 & 215 & 2.856 & 87 & 3.299 \\
H5 & 208 & 2.750 & 94 & 3.234 \\
H6 & 212 & 2.972 & 98 & 3.255 \\
H7 & 213 & 3.014 & 101 & 3.317 \\
H8 & 212 & 2.991 & 93 & 3.204 \\
H9 & 211 & 3.085 & 99 & 3.303 \\
M1 & 420 & 2.305 & 112 & 2.955 \\
M2 & 377 & 2.626 & 137 & 3.117 \\
M3 & 414 & 2.935 & 198 & 3.192 \\
M4 & 432 & 2.965 & 208 & 3.216 \\
M5 & 417 & 2.609 & 165 & 3.085 \\
M6 & 421 & 2.888 & 196 & 3.204 \\
N1 & 209 & 1.526 & 28 & 2.750 \\
N2 & 75 & 3.013 & 50 & 3.300 \\
\midrule
SD across the 17 non-null cells (high-Y column) & --- & --- & --- & 0.102 \\
SD across all 19 cells (incl. N1, N2; high-Y column) & --- & --- & --- & 0.142 \\
\bottomrule
\end{tabular}
\end{table}


\subsection{Capability preservation}
\label{sec:discussion_capability}

\begin{table}[t]
\centering
\small
\caption{Capability preservation. For each path, we restrict to samples on which the strongest mono baseline (Path 1: M4, Path 2: M3, Path 3: M4) achieves valid closure. Each cell reports the fraction of those same samples on which the row-cell also achieves valid closure. Values near $1.0$ indicate the cell does not sacrifice mono strengths. \textbf{Bold} marks preservation $\geq 0.90$ (well-preserved); \emph{italic} marks preservation $< 0.50$ (materially collapsed).}
\label{tab:capability_preservation}
\begin{tabular}{lrrr}
\toprule
Cell & Path 1 preservation & Path 2 preservation & Path 3 preservation \\
\midrule
M1 & 0.595 & \emph{0.422} & \emph{0.440} \\
M2 & \emph{0.491} & 0.711 & 0.615 \\
M3 & 0.741 & \textbf{1.000} & 0.835 \\
M4 & \textbf{1.000} & 0.889 & \textbf{1.000} \\
M5 & 0.767 & 0.689 & 0.606 \\
M6 & 0.802 & 0.844 & 0.798 \\
\midrule
H1 & 0.778 & \textbf{0.909} & 0.758 \\
H2 & \emph{0.492} & 0.761 & 0.636 \\
H3 & 0.762 & 0.804 & \emph{0.455} \\
H4 & 0.825 & 0.565 & 0.782 \\
H5 & 0.730 & 0.891 & 0.527 \\
H6 & 0.762 & 0.826 & 0.745 \\
H7 & 0.778 & \textbf{0.913} & 0.800 \\
H8 & 0.794 & 0.804 & 0.818 \\
H9 & 0.762 & 0.891 & 0.873 \\
H10 & \emph{0.349} & 0.739 & 0.709 \\
H11 & \textbf{0.905} & 0.870 & 0.800 \\
\midrule
N1 & \emph{0.333} & \emph{0.152} & \emph{0.291} \\
N2 & \emph{0.000} & 0.891 & \emph{0.000} \\
N3 & \emph{0.000} & \emph{0.000} & \emph{0.000} \\
\bottomrule
\end{tabular}
\end{table}

Table~\ref{tab:capability_preservation} reports the capability-preservation rate: for each path, the fraction of samples on which the strongest mono baseline (M4 for Paths~1 and~3, M3 for Path~2) achieves valid closure, restricted to those on which each other cell also achieves closure. Well-composed hetero cells --- H11, H7, H9, H1 --- preserve 76--91\% of reference capability across the three paths (per-path minimum $0.758$, H1 on Path~3; maximum $0.913$, H7 on Path~2). Two cells sacrifice materially:

\begin{itemize}
  \item H10 preserves only $34.9\%$ of Path~1 reference capability. This is another manifestation of the phi-4-in-test-stage pathology: samples on which M4 (gemma-4 mono) achieves valid closure are precisely the samples on which H10 has phi-4 fail to produce filter-passing tests.
  \item H4 preserves only $56.5\%$ of Path~2 capability. Placing llama3.2:3B at $L_{\textsf{code}}$ costs the pipeline $44\%$ of reference-pass-rate success on samples that qwen3.6 mono can synthesise cleanly. The heterogeneity gain of $\Delta = 0.285$ on Path~1 observed in the stage-attribution decomposition (Section~\ref{sec:results_stage_attribution}) hides the fact that H4 pays this capability cost on Path~2 on the samples the strongest mono baseline handles well.
\end{itemize}

This analysis quantifies a design tension: heterogeneity offers gains via specialization but incurs losses when a stage's model is inadequate for that stage. The design implication is that heterogeneous pipelines are not uniformly ``better'' than the strongest mono baseline; they should be selected based on the target sample distribution and the availability of a strong-enough model at each stage.

\subsection{Threshold sensitivity}
\label{sec:discussion_threshold_sensitivity}

The RQ1 conclusion of Section~\ref{sec:discussion_rq1} holds at every $(\tau, \rho) \in \{0, 0.5\} \times \{2, 3\}$ on Path~1: no mono cell dominates all hetero cells at any threshold pair in the operational band, and Kendall's $\tau_{\textsf{rank}}$ between the reference threshold and every alternative exceeds $0.65$. Path~3 sensitivity is markedly weaker: at $\tau = 0.4$ (a strict BERTScore threshold), Kendall's $\tau_{\textsf{rank}}$ falls to $-0.15$ (non-significant), consistent with the RQ2 finding that BERTScore is not carrying a semantic-equivalence signal on Path~3. Full $9$-cell threshold grid and the conditional-on-metric denominator convention $V'$ appear in Appendix~\ref{sec:appendix_threshold}.

\subsection{Design guidelines}

We distil three design guidelines for practitioners assembling multi-SLM SE pipelines.

\textbf{G1: Place the strongest available model at the bottleneck stage.} The stage-attribution decomposition (Section~\ref{sec:results_stage_attribution}) and per-mutation-operator table (Section~\ref{sec:results_per_operator}) both identify per-stage strengths that dominate closure success. When choosing a stage assignment, the correct heuristic is not ``pick the strongest model overall'' but ``identify which stage the pipeline is bottlenecked at, and place the strongest model there.'' A per-operator capability decomposition (phi-4 on predicate reasoning; qwen-coder on comparison operators) suffices as an initial oracle; empirical validation via mutation testing sharpens it.

\textbf{G2: Match test-generator to the docstring style upstream, not to the model's general capability.} The phi-4-in-test-stage pathology (Section~\ref{sec:discussion_rq3}) shows that a model's capability at generating tests \emph{de novo} does not transfer to generating tests conditioned on a docstring produced by a different model. When phi-4 is downstream of a qwen3.6 or mistral docstring, its test-generation quality collapses. The practical implication is that stage models should be chosen jointly, not marginally.

\textbf{G3: Use two-signal validity when reporting closure rates.} The judge--metric decomposition (Section~\ref{sec:discussion_rq2}) shows that any single automated metric can fail structurally --- BERTScore on Path~3 explains at most $4.5\%$ of variance in the judge rating within either benchmark and shows a categorical benchmark-conditional disagreement flip. Reporting closure rate as ``metric $> \tau$'' alone risks systematic misestimation. Algorithm~\ref{alg:validity}'s two-signal strict-AND policy sacrifices some closure counts but produces a validity signal whose disagreement modes (Table~\ref{tab:disagreement_by_path}) are interpretable and can be reported honestly.
